\scanpage{222}

\mspage{96}

Comme par hypothèse $f=\beta_i\alpha_i$ est dans $W$, et $\alpha_i\in W$, il s'ensuit par 7.1.4. a) que $\beta_i\in W$. Il faudra que nous sachions de plus que $\beta_i$ est mono, donc $\beta_i\in\mathrm{Mon}\cap W=C$, ce qui achèvera \deleted{de vérifier} d'établir la condition manquante b) de 7.3. C'est la condition \deleted{(7.1.1.)}\inserted{d'exactitude} qu'il me faut visiblement rajouter aux hypothèses \deleted{du} 7.1. 3°). Ayant fait cela, le lemme 7.3. s'applique et établit 7.1.3°.

\uline{Démonstration de} 7.2. Posons

\[
C=\mathrm{Mon},\quad\mathrm{TC}=C\cap W,\quad\mathrm{TC}_0=\mathrm{TC}^{\Pi},
\]

on a donc par 7.1. 3°, on a

\[
\underbrace{\mathrm{TC}_0}_{\substack{\text{ess}^{\mathrm{t}}\\\text{petite}}}\subset\mathrm{TC}\subset\widetilde{\mathrm{TC}_0}=\widetilde{\mathrm{TC}}
\]

et par le th. 1 \S1 (p. \deleted{\unsure{10}}12), cor. \unsure{1|2}, pour vérifier que $\mathrm{TC}=\widetilde{\mathrm{TC}_0}$ (i.e $\mathrm{TC}$ est $Q$-clos \unsure{à|:} g.) il \deleted{\unsure{sert}} \unsure{revient} au $=$ de voir qu'il est Q-stable, i.e \unsure{S|S°.}) être stable par \inserted{compositions, par} cochgt \deleted{de} \unsure{base|b.}, par \uline{$\lim$} ordinales strictes, par fact. directe. Or les deux dernières sont induites par le fait que $\inserted{\mathrm{TC}=}C\cap W$ est stable par \uline{$\lim$} filtrantes, vu que $W$ l'est, et $C$ aussi par la condition \textcircled{e} \inserted{(p. 76)} sur $M$. D'autre part, $C=\mathrm{Mon}$ et $W$ étant stables par composition, idem pour $C\cap W$. Donc il reste le cochgt. de
% archive sheet number: 221
